\documentclass{article}
\usepackage[utf8]{inputenc}
\usepackage{amsmath}
\usepackage{amssymb}
\usepackage{geometry}
\geometry{a4paper, margin=1in}

\title{Foundations of Reinforcement Learning: MDPs and Q-Learning}
\author{Reinforcement Learning Primer}
\date{\today}

\begin{document}

\maketitle

\section{The Markov Decision Process (MDP) Framework}
The control problem in reinforcement learning is formally defined using a Markov Decision Process, represented by the tuple $(\mathcal{S}, \mathcal{A}, P, R, \gamma)$.

\subsection{Components}
\begin{itemize}
    \item \textbf{State Space ($\mathcal{S}$):} The set of all possible environment configurations.
    \item \textbf{Action Space ($\mathcal{A}$):} The set of all available moves for the agent.
    \item \textbf{Transition Probability ($P$):} The mapping $P(s' | s, a)$, defining the probability of reaching state $s'$ from state $s$ via action $a$.
    \item \textbf{Reward Function ($R$):} The immediate feedback $R(s, a)$ received after taking an action.
    \item \textbf{Discount Factor ($\gamma$):} A constant $\gamma \in [0, 1)$ that balances immediate and future rewards.
\end{itemize}

\subsection{The Objective Function}
The agent's goal is to maximize the expected discounted return $G_t$:
\begin{equation}
    G_t = \sum_{k=0}^{\infty} \gamma^k r_{t+k+1}
\end{equation}

\section{The Q-Learning Algorithm}
Q-learning is a model-free, off-policy algorithm used to find the optimal action-value function $Q^*(s, a)$.

\subsection{The Bellman Optimality Equation}
The optimal Q-value, denoted by $Q^*(s, a)$, is the maximum expected return achievable from a given state-action pair $(s, a)$, and it satisfies the Bellman optimality equation:
\begin{equation}
    Q^*(s, a) = \mathbb{E} \left[ r + \gamma \max_{a'} Q^*(s', a') \mid s, a \right], 
    \forall (s, a) \in \mathcal{S} \times \mathcal{A}
\end{equation}
where $s'$ is the next state resulting from taking action $a$ in state $s$.
We denote 
\begin{equation}
    \left[ r + \gamma\max_{a'} Q^*(s', a') \mid s, a \right] = F(Q^*, s, a, \omega),
    \forall (s, a) \in \mathcal{S} \times \mathcal{A}
\end{equation}
where $\omega$ represents the randomness in the environment. Then, the 
optimal Q-value can be expressed as fixed point of the stochastic Bellman operator $F[Q]$:
\begin{equation}
    Q^*(s, a) = \mathbb{E} \left[ F(Q^*, s, a, \omega) \right], \forall (s, a) \in \mathcal{S} \times \mathcal{A}.
\end{equation}

\subsection{Temporal Difference (TD) Update Rule}
It is well-known in the literature of Stochastic Approximation that 
the fixed point of $x = \mathbb E [f(x, \omega)]$ can be found 
by the following iterative procedure: For an appropriate condition 
of $f$ and a sequence of i.i.d. samples $\{\omega_n\}$, 
the sequence $\{x_n\}$ defined by
\begin{equation}
    x_{n+1} = (1 - \alpha) x_n + \alpha f(x_n, \omega_n)
\end{equation}
where $\alpha$ is the learning rate. Applying this result to
the Bellman optimality equation, we can derive the Q-learning update rule. 
\begin{equation}
    Q(s, a) \leftarrow Q(s, a) + \alpha \left[ r + \gamma \max_{a'} Q(s', a') - Q(s, a) \right]
\end{equation}
where $\alpha$ is the \textbf{learning rate}.

\section{Control and Policy}
Once the Q-values converge to $Q^*$, the optimal policy $\pi^*$ is derived by selecting the action with the highest value in each state:
\begin{equation}
    \pi^*(s) = \arg\max_{a \in \mathcal{A}} Q^*(s, a)
\end{equation}
The above Q-learning algorithm is pure exploitation or zero-greedy. 
In practice, we often use an $\epsilon$-greedy policy to
balance exploration and exploitation during learning. 
The $\epsilon$-greedy policy selects a random action with probability $\epsilon$ and the action with the highest Q-value with probability $1 - \epsilon$.

\end{document}